Т% Separate concise edition; the full draft and its assets are preserved.
% Official Cambridge PASA template (2022); vendor class is unmodified.
% pas.cls embeds a pre-2020 \document implementation. Keep the current kernel
% entry point so modern package hooks run (TeX Live 2025 / current Overleaf).
\let\PasaKernelDocument\document
% The class loads cleveref itself; amsmath must precede it.
\RequirePackage{amsmath}
\documentclass[nodoi,nocopyright]{pas}
\let\document\PasaKernelDocument
% Use the publisher's page size in modern hyperref as well as in pdfTeX.
\setlength{\paperwidth}{\trimwidth}
\setlength{\paperheight}{\trimheight}
% The template prints a publisher copyright even with its nocopyright option.
% Suppress that unassigned publication metadata for this author manuscript.
\makeatletter
\patchcmd{\ps@myplain}{\copyright\ \@cpr}{}{}{}
\patchcmd{\ps@myplain}{\copyright\ \@cpr}{}{}{}
\makeatother

\usepackage[T1]{fontenc}
\usepackage[utf8]{inputenc}
\usepackage[british]{babel}
\usepackage{array}
\usepackage{microtype}
\usepackage{placeins}
% Permit full-width floats at page bottoms in the two-column manuscript.
\usepackage{stfloats}
\usepackage{capt-of}
\hypersetup{
  hidelinks,
  pdftitle={JWST/NIRCam F090W and F150W Surface Brightness Fluctuation Distances},
  pdfauthor={I. A. Bukanov}
}

\newcommand{\mbar}{\overline{m}}
\newcommand{\Mbar}{\overline{M}}
% Fit wide tables downwards only; never enlarge their native PASA font.
\newsavebox{\papertablebox}
\newcommand{\papertable}[2]{%
  \sbox{\papertablebox}{\tablefont\setlength{\tabcolsep}{4pt}#2}%
  \ifdim\wd\papertablebox>#1\relax
    \resizebox{#1}{!}{\usebox{\papertablebox}}%
  \else\usebox{\papertablebox}\fi
}
\newcommand{\paperfigure}[2]{%
  \IfFileExists{#1}{%
    \includegraphics[width=#2]{#1}%
  }{%
    \fbox{\parbox[c][42mm][c]{0.90\linewidth}{\centering
      Figure not generated.\\
      Run the corresponding analysis-notebook cell.\\[1mm]
      \footnotesize\nolinkurl{#1}}}%
  }%
}

\begin{document}
\lefttitle{Publications of the Astronomical Society of Australia}
\righttitle{I. A. Bukanov}
\jnlPage{1}{\pageref{LastArticlePage}}
\jnlDoiYr{Draft}
\renewcommand{\doitext}{}
\articletitt{Research Paper}
\title{JWST/NIRCam F090W and F150W Surface Brightness Fluctuation Distances
from GO-3055: TRGB-Anchored Calibrations and Internal Accuracy Tests}
\author{\gn{I. A.} \sn{Bukanov}$^1$}
\affil{$^1$Lomonosov Moscow State University, Faculty of Space Research, Moscow, Russia}
% Submission metadata remain for the authors to confirm, not for the renderer
% to invent: co-authors, corresponding-author email, funding, conflicts,
% acknowledgements, and the required AI-assistance disclosure.
\corresp{I. A. Bukanov (contact details to be confirmed)}
\citeauth{Bukanov, I. A. Draft manuscript; author list to be agreed.}

\begin{abstract}
We measure surface brightness fluctuations (SBF) in JWST/NIRCam F090W and
F150W images of 14 early-type galaxies from GO-3055, using a common
measurement procedure and individual tip of the red giant branch (TRGB)
distances. All galaxies enter the primary calibration. F150W does not
require a colour correction in this sample: its adopted constant absolute
magnitude is $-3.208\pm0.024_{\rm stat}\pm0.051_{\rm common}$ mag, with
intrinsic scatter 0.081 mag. Adding a linear colour term worsens the
leave-one-out (LOO) root-mean-square (RMS) distance-modulus residual
from 0.0925 to 0.0981 mag.
In F090W, the relation
$\Mbar_{090,0}=-1.9617+1.8890[(F090W-F150W)_0-0.5703]$ mag reduces
the LOO RMS from 0.0992 to 0.0684 mag and has intrinsic scatter 0.0489 mag.
The median full absolute distance uncertainties are 0.0984 and 0.0752 mag
for F150W and F090W, respectively, including the common 0.047-mag TRGB scale.
We separate image-measurement, population, calibration and common-scale
uncertainties, and compare distances with HST F110W and F160W results on
a consistent error convention. These are internal cross-validations, not
independent external accuracy tests. Correlated image noise, empirical
PSF normalisation and population-dependent systematics remain limitations.
\end{abstract}
\begin{keywords}
Galaxy distances -- Elliptical galaxies -- Infrared photometry
\end{keywords}
\maketitle\endgroup

\section{Introduction}

Surface brightness fluctuations measure the variance of unresolved stellar
light per unit mean flux \citep{TonrySchneider1988}. Their apparent amplitude
decreases with distance, while their absolute magnitude depends on the
stellar luminosity function. Near-infrared SBF benefits from luminous evolved
stars but can require a stellar-population correction
\citep{Jensen2003,Jensen2015,Jensen2021}. Colour is a potential predictor of
this correction, not an ingredient of the image-based SBF measurement.

The TRGB--SBF project transfers geometrically anchored tip of the red giant
branch (TRGB) distances to SBF. \citet[Paper III]{Jensen2025} recalibrated
Hubble Space Telescope (HST) F110W SBF using JWST TRGB distances and mean Virgo/Fornax
anchors. \citet[Paper IV]{Chazov2026} subsequently provided individual
JWST TRGB distances. We use these anchors to measure and calibrate SBF
directly in both JWST F090W and F150W. The two bands are processed with
one methodology, allowing a controlled comparison of their colour response
and internal distance precision.

We compare constant and linear calibrations before choosing the distance
estimator. The primary experiment retains all 14 galaxies, without manual
outlier removal. Technical derivations, complete residual-image montages and
extended sensitivity tests are collected in the accompanying
\emph{Technical notes}, retained alongside this concise manuscript;
the evidence needed to interpret the adopted distances is given here.

\section{Data and common measurement method}
\label{sec:method}

\subsection{Images, distance anchors and residuals}

The input data are public GO-3055 James Webb Space Telescope (JWST)
Near Infrared Camera (NIRCam) i2d mosaics in MJy~sr$^{-1}$,
with native pixel scale approximately $0.03123$ arcsec. The sample comprises
three Fornax galaxies, nine in the broad Virgo region and two other nearby
group galaxies. The Virgo-region label is not strict cluster membership;
in particular it includes NGC 4636 and foreground NGC 4697.

We use the individual Paper IV thin-slice TRGB moduli throughout the primary
analysis. Thin slices were fixed before inspecting the present calibration
residuals; Paper IV's preferred thick-slice distances are retained as a
sensitivity test. The SBF regions are circular annuli at
$8.2$--$16.4$ and $16.4$--$32.8$ arcsec, following Paper III.
The outer radius spans 1.65--3.17 kpc across the sample, so these are not
matched physical apertures. The auxiliary elliptical-annulus test of Jensen
et al. is not reproduced.

In each band we subtract a scalar sky, fit a smooth isophotal galaxy model
with free ellipticity and position angle, and mask invalid pixels and
compact sources. Model intensities and geometry are interpolated between
isophotes; filled detector gaps are never used for SBF. For sky-subtracted
image $I_i$ and positive galaxy model $M_i$, define
\begin{equation}
 Q_i=(I_i-M_i)/\sqrt{M_i}.
 \label{eq:normalisation}
\end{equation}
This removes the leading dependence of stellar variance on local galaxy
light. A single median and scale are estimated with iterative clipped
statistics over the complete catalogue-unmasked, finite, positive-model
support, not separately in each annulus. Values of $Q_i$ beyond
median $\pm3.5$ scale are replaced by the boundary values.
This is \emph{winsorisation}, not pixel rejection, and is an explicit
departure from Jensen's procedure.

The full normalised, winsorised field is saved before selecting annuli.
For binary annular window $W_i$, the Fourier input is
\begin{equation}
 R_i=W_i\left(Q_i^{\rm win}-\langle Q^{\rm win}\rangle_W\right).
 \label{eq:fft-input}
\end{equation}
The mean is removed within each usable ring; masked pixels are absent
rather than measurements of zero flux. Complete F150W and F090W montages
in the Technical notes show the actual measurement support, black masked
regions and the three annular boundaries.

\subsection{Fourier amplitude and apparent SBF magnitude}

The measured two-dimensional power spectrum is
$P(\boldsymbol{k})=|\mathcal F[R]|^2/N_{\rm use}$, where $N_{\rm use}$
is the number of usable annular pixels.
We average it into 79 radial bins defined by 80 edges. Bin errors are
standard errors of Fourier-mode power, counting each conjugate pair of a
real image once. The window can still correlate distinct modes; this is not
a complete spectral-covariance treatment.

For each of five time- and detector-position-matched STPSF models
\citep{Perrin2014}, the unit-flux point-spread function (PSF) is convolved with 64 independent
unit white-noise fields. Applying the same mask, annular mean subtraction
and Fourier normalisation gives an expectation spectrum $E_j(k)$.
The adopted PSF stamps are 129 pixels across. Gaussian simulation fields
probe the second-order PSF/window response; they do not require a Gaussian
one-point histogram of the observed residual image.

For each PSF we fit
\begin{equation}
 P(k)=P_{0,j}E_j(k)+P_{1,j}
 \label{eq:power-model}
\end{equation}
by inverse-bin-variance weighted least squares (WLS). The fit interval is
$0.04\leq k\leq0.25$ pixel$^{-1}$, with a lower bound of ten modes across
the shorter crop axis. The coefficient covariance
$(X^{\mathsf T}WX)^{-1}$ is multiplied by $\max(\chi^2_\nu,1)$.
The adopted amplitudes are medians over the five PSFs; windows beginning
at $k_{\min}=0.01$ and 0.03 test low-frequency sensitivity.

The fluctuation power is $P_{\rm SBF}=P_0-P_r$, where $P_r$ is the
unresolved globular-cluster and background-galaxy contribution. It is
estimated by flux-squared integration of a Gaussian globular-cluster
luminosity function and a power-law background population below the source
catalogue limit. The adopted globular-cluster width is 1.2 mag.
Since Equation~\ref{eq:normalisation} has already normalised the variance
by the mean galaxy light, $P_{\rm SBF}$ is a fluctuation flux: it must not
be divided by the model a second time. For pixel solid angle
$\Omega_{\rm pix}$,
\begin{equation}
 \mbar_X=-2.5\log_{10}(P_0-P_r)
 -2.5\log_{10}\left(\frac{10^6\Omega_{\rm pix}}{3631}\right).
 \label{eq:mbar}
\end{equation}

The two annular magnitudes are combined with their saved normalised
inverse-spectral-variance weights, not area weights. F090W is registered
and PSF-matched to F150W for the matched-annulus colour measurement.
Each band's combined colour uses that band's SBF weights; the two combined
colours can therefore differ slightly. Annular magnitude and colour
differences remain radial diagnostics, not random-error floors.

\subsection{Extinction and calibration}
\label{sec:calibration-method}

We use Paper IV's $A_{090}$ with $E(B-V)=A_{090}/1.4156$ and
$A_{150}=0.6021E(B-V)$. Then
$\mbar_{X,0}=\mbar_{X,\rm obs}-A_X$ and
$C_0=C_{\rm obs}-(A_{090}-A_{150})$.
The absolute calibration points are
$\Mbar_{X,0}=\mbar_{X,0}-\mu_{\rm TRGB}$.
Colour first enters at this stage, through a trial calibration $f_X(C_0)$;
it never enters Equation~\ref{eq:power-model}.

Constant and linear relations are fitted by a Gaussian likelihood with a
free non-negative intrinsic scatter $\sigma_{\rm int}$. For effective
variance
\begin{equation}
 V_i=\sigma^2_{\Mbar,i}+[f'(C_i)\sigma_{C,i}]^2
 -2f'(C_i)\operatorname{Cov}(C_i,\Mbar_i)+\sigma_{\rm int}^2,
 \label{eq:variance-fit}
\end{equation}
the objective contains both residual-squared/$V_i$ and $\ln V_i$.
Model selection uses AICc and leave-one-galaxy-out (LOO) prediction.
Each target is omitted before calibrating its distance.
Quadratic, broken-line, logarithmic and exponential laws remain exploratory;
they do not replace the primary models solely because a curve looks tighter.

\section{Common uncertainty model}
\label{sec:errors}

An apparent SBF measurement, an absolute calibration point and a predicted
distance have different uncertainties. Shared scale terms are correlated
between galaxies and are not independent regression weights.

\subsection{Image measurement and ring combination}

Writing $h=2.5/\ln10$ and MAD for the median absolute deviation, the
ring-level spectral error is
\begin{equation}
 \sigma_{{\rm PS},r}=\max\left[
 \frac{h\,\sigma_{P_0,\rm WLS}}{P_0-P_r},
 1.4826\,{\rm MAD}\{\mbar_r(k_{\min})\}\right].
 \label{eq:spectrum-error}
\end{equation}
The maximum treats two estimates of uncertainty in the same amplitude,
rather than adding them as independent physical sources. It is a
conservative implementation choice, not Jensen's exact error prescription.
Other ring-level contributions are
\begin{align}
 \sigma_{{\rm PSF},r}&=\frac{h\,1.4826\,{\rm MAD}\{P_{0,r,j}\}}{P_{0,r}-P_{r,r}},
 \\
 \sigma_{{\rm sky},r}&=\frac{h\,s_{\rm sky}}{\langle M\rangle_r},
 \qquad
 \sigma_{P_r,r}=\frac{h(0.25)P_{r,r}}{P_{0,r}-P_{r,r}}.
\end{align}
PSF scatter is measured in power before propagation to magnitudes.
Here $s_{\rm sky}$ is the scatter of independent sky estimates.
The sky term is a first-order normalisation sensitivity, not a complete
sky/model refit. The $0.25P_r$ allowance is assumed, rather than derived
from artificial-source completeness tests.

For saved weights $w_r=\sigma_r^{-2}/\sum_s\sigma_s^{-2}$, independent
spectral terms combine as $\sum_r w_r^2\sigma_{{\rm PS},r}^2$.
Each PSF, sky and $P_r$ term is combined linearly across rings,
$\sigma_j=\sum_r w_r\sigma_{j,r}$, because its shift is treated as coherent.
The internal image-measurement variance is
\begin{equation}
 \sigma_{\mbar,\rm int}^2
 =\sigma_{\rm PS}^2+\sigma_{\rm PSF}^2+
  \sigma_{\rm sky}^2+\sigma_{P_r}^2.
 \label{eq:measurement-error}
\end{equation}
No annular half-difference is added to this quantity.

\subsection{Calibration-point and distance errors}

The adopted reddening uncertainty is
$\sigma_{A_{090}}^2=(0.16A_{090})^2+0.005^2$.
The individual Paper IV error already contains $A_{090}$; we remove that
term before restoring the appropriate shared-extinction covariance.
Defining $\sigma_{\mu,\rm nr}^2=\sigma_{\mu,\rm IV}^2-\sigma_{A_{090}}^2$
and $\sigma_{\Delta A}=(1.4156-0.6021)\sigma_{E(B-V)}$, we obtain
\begin{align}
 \sigma_{\Mbar_{150}}^2
 &=\sigma_{\mbar_{150},\rm int}^2+\sigma_{\mu,\rm nr}^2+\sigma_{\Delta A}^2,
 \\
 \sigma_{\Mbar_{090}}^2
 &=\sigma_{\mbar_{090},\rm int}^2+\sigma_{\mu,\rm nr}^2.
 \label{eq:point-errors}
\end{align}
The foreground colour--magnitude covariance is
$-\sigma_{\Delta A}^2$ in F150W and zero in F090W, where the same
$A_{090}$ cancels against the TRGB anchor. These statements concern only
foreground extinction, not every shared-image covariance.

The formal photometric colour error comes from the pixel-colour standard
error within matched rings, combined as $\sum_r w_r^2\sigma_{C,r}^2$.
Spatially correlated pixels and image-wide photometric systematics are
not fully represented. A separate 0.01-mag colour-floor sensitivity test
does not reverse the adopted model choice in either band.

With $R_X=A_X/E(B-V)$ and $\sigma_E=\sigma_{E(B-V)}$, a target has
$\mu=\mbar_{X,\rm obs}-R_XE(B-V)-f_X(C_0)$. Its measurement
and reddening propagation is
\begin{align}
 \sigma_{\rm target}^2={}&\sigma_{\mbar_X,\rm int}^2+
 [f'_X\sigma_{C,\rm phot}]^2 \nonumber\\
 &+[R_X-f'_X(R_{090}-R_{150})]^2\sigma_E^2.
 \label{eq:target}
\end{align}
Thus shared reddening is counted once. The predictive error is
\begin{align}
 \sigma_{\mu,\rm int}^2&=\sigma_{\rm target}^2+
 \sigma_{\rm int}^2+\sigma_{\rm cal}^2,\\
 \sigma_{\mu,\rm abs}^2&=\sigma_{\mu,\rm int}^2+0.047^2.
 \label{eq:distance-errors}
\end{align}
The scatter and finite-calibration terms are determined from each
13-galaxy LOO training set; $\sigma_{\rm cal}$ is the standard deviation
of bootstrap predictions (300 realisations for F150W, 350 for F090W).
The 0.047-mag common TRGB/NGC 4258 scale \citep{Reid2019,Chazov2026} moves
all absolute distances together.

For the absolute F150W zero point, common TRGB, flux-scale and PSF-stamp
terms are 0.047, 0.0108 and 0.0167 mag, combining to 0.0510 mag.
The latter two cancel to first order between apparent and calibrated
absolute magnitudes in a same-filter, same-pipeline distance.
They are not added again as independent distance errors.
We convert moduli using
$D=10^{(\mu-25)/5}$ Mpc and
$\sigma_D\simeq(\ln10/5)D\sigma_\mu$.
Validation residuals $\mu_{\rm SBF}-\mu_{\rm TRGB}$ instead include the
individual target anchor error with its reddening covariance; the common
TRGB zero point cancels. This validation uncertainty, the empirical LOO
RMS and the full absolute distance error must not be interchanged.

\section{Results in F150W and F090W}
\label{sec:results}

\subsection{Measurements and model selection}

Tables~\ref{tab:f150-measurements} and \ref{tab:f090-measurements} retain
all 14 objects. The tabulated absolute magnitudes are image measurements
minus individual TRGB moduli, not values predicted by a colour law.
Figure~\ref{fig:calibrations} and Table~\ref{tab:models} compare both
primary models in each band.

\begin{table*}[t]
\centering
\caption{F150W extinction-corrected measurements. $C_{090-150,0}$ is the
matched colour; $\mbar_{150,0}$ is apparent SBF and
$\Mbar_{150,0}=\mbar_{150,0}-\mu_{\rm TRGB}$. The individual
$\sigma_{\Mbar}$ follows Equation~\ref{eq:point-errors}, excluding common
zero points. $|\Delta\mbar|_{\rm ann}$ is the full absolute annular
difference, not an error term; $\Delta\mu_{\rm LOO,const}$ is the
colour-independent validation residual. HQ IV denotes Paper IV's quality
flag and does not select the primary sample. All magnitudes are in mag.}
\label{tab:f150-measurements}
\papertable{\textwidth}{\begin{tabular}{llrrrrrrrc}
\toprule
Galaxy & Environment & $C_{090-150,0}$ & $\overline m_{150,0}$ &
$\overline M_{150,0}$ & $\sigma_{\overline M}$ &
$|\Delta\overline m|_{\rm ann}$ & $\mu_{\rm TRGB}$ &
$\Delta\mu_{\rm LOO,const}$ & HQ IV \\
 & & (mag) & (mag) & (mag) & (mag) & (mag) & (mag) & (mag) & \\
\midrule
NGC 1380 & Fornax & 0.510 & 28.087 & -3.321 & 0.029 & 0.099 & 31.408 & -0.122 & $+$ \\
NGC 1399 & Fornax & 0.570 & 28.157 & -3.341 & 0.034 & 0.047 & 31.498 & -0.143 & $-$ \\
NGC 1404 & Fornax & 0.559 & 28.177 & -3.189 & 0.028 & 0.110 & 31.366 & +0.020 & $+$ \\
NGC 1549 & Other & 0.533 & 27.930 & -3.301 & 0.029 & 0.133 & 31.231 & -0.101 & $-$ \\
NGC 3379 & Other & 0.557 & 27.012 & -3.066 & 0.030 & 0.269 & 30.078 & +0.152 & $+$ \\
NGC 4374 & Virgo region & 0.566 & 27.947 & -3.165 & 0.031 & 0.094 & 31.112 & +0.046 & $+$ \\
NGC 4406 & Virgo region & 0.616 & 27.849 & -3.247 & 0.030 & 0.097 & 31.096 & -0.042 & $+$ \\
NGC 4472 & Virgo region & 0.625 & 27.713 & -3.303 & 0.035 & 0.029 & 31.016 & -0.103 & $-$ \\
NGC 4486 & Virgo region & 0.652 & 27.808 & -3.238 & 0.055 & 0.054 & 31.046 & -0.033 & $-$ \\
NGC 4552 & Virgo region & 0.569 & 27.850 & -3.073 & 0.030 & 0.175 & 30.923 & +0.146 & $+$ \\
NGC 4621 & Virgo region & 0.565 & 27.688 & -3.133 & 0.028 & 0.173 & 30.821 & +0.081 & $+$ \\
NGC 4636 & Virgo region & 0.604 & 27.890 & -3.180 & 0.031 & 0.110 & 31.070 & +0.031 & $+$ \\
NGC 4649 & Virgo region & 0.623 & 27.873 & -3.136 & 0.030 & 0.043 & 31.009 & +0.077 & $-$ \\
NGC 4697 & Virgo region & 0.512 & 27.095 & -3.229 & 0.032 & 0.173 & 30.324 & -0.023 & $+$ \\
\bottomrule
\end{tabular}
}
\end{table*}

\begin{table*}[t]
\centering
\caption{F090W extinction-corrected measurements, with the same object
selection and definitions as Table~\ref{tab:f150-measurements}.
$\sigma_{\mbar}$ is the internal image-measurement error from
Equation~\ref{eq:measurement-error}; $\sigma_{\Mbar}$ additionally contains
the individual TRGB term with the shared $A_{090}$ cancellation.
The colour uses the F090W annular weights, rather than the F150W weights.}
\label{tab:f090-measurements}
\papertable{\textwidth}{\begin{tabular}{llrrrrrrc}
\toprule
Galaxy & Env. & $C_{090-150,0}$ & $\overline m_{090,0}$ & $\sigma_{\overline m}$ & $\overline M_{090,0}$ & $\sigma_{\overline M}$ & $|\Delta\overline m|_{\rm ann}$ & HQ IV \\
 & & (mag) & (mag) & (mag) & (mag) & (mag) & (mag) & \\
\midrule
NGC~1380 & Fornax & 0.511 & 29.328 & 0.020 & -2.080 & 0.032 & 0.039 & $+$ \\
NGC~1399 & Fornax & 0.569 & 29.616 & 0.023 & -1.882 & 0.034 & 0.079 & $-$ \\
NGC~1404 & Fornax & 0.565 & 29.414 & 0.016 & -1.952 & 0.029 & 0.006 & $+$ \\
NGC~1549 & Other & 0.539 & 29.214 & 0.027 & -2.017 & 0.037 & 0.045 & $-$ \\
NGC~3379 & Other & 0.570 & 28.156 & 0.017 & -1.922 & 0.032 & 0.156 & $+$ \\
NGC~4374 & Virgo region & 0.566 & 29.163 & 0.023 & -1.949 & 0.034 & 0.029 & $+$ \\
NGC~4406 & Virgo region & 0.615 & 29.130 & 0.024 & -1.966 & 0.034 & 0.040 & $+$ \\
NGC~4472 & Virgo region & 0.626 & 29.128 & 0.016 & -1.888 & 0.031 & 0.018 & $-$ \\
NGC~4486 & Virgo region & 0.652 & 29.199 & 0.057 & -1.847 & 0.062 & 0.056 & $-$ \\
NGC~4552 & Virgo region & 0.574 & 29.025 & 0.019 & -1.898 & 0.033 & 0.088 & $+$ \\
NGC~4621 & Virgo region & 0.571 & 28.806 & 0.019 & -2.015 & 0.031 & 0.075 & $+$ \\
NGC~4636 & Virgo region & 0.603 & 29.097 & 0.024 & -1.973 & 0.034 & 0.041 & $+$ \\
NGC~4649 & Virgo region & 0.622 & 29.254 & 0.013 & -1.755 & 0.027 & 0.045 & $-$ \\
NGC~4697 & Virgo region & 0.517 & 28.198 & 0.016 & -2.126 & 0.031 & 0.057 & $+$ \\
\bottomrule
\end{tabular}
}
\end{table*}

\begin{table*}[t]
\centering
\caption{Constant and linear calibrations for the same 14 galaxies.
The linear form is $a+b(C_0-C_p)$, with $C_p=0.5675$ mag for F150W
and 0.5703 mag for F090W. Lower AICc and LOO RMS are preferred;
the adopted models are marked by an asterisk.
$a$, $\sigma_{\rm int}$ and LOO RMS are in mag.}
\label{tab:models}
\papertable{\textwidth}{%
\begin{tabular}{llrrrrr}
\toprule
Band & Model & $a$ & $b$ & $\sigma_{\rm int}$ & AICc & LOO RMS\\
\midrule
F150W & Constant$^*$ & -3.2078 & 0 & 0.0805 & -23.75 & 0.0925\\
F150W & Linear & -3.2094 & 0.2315 & 0.0797 & -20.61 & 0.0981\\
F090W & Constant & -1.949 & 0 & 0.086 & -21.84 & 0.0992\\
F090W & Linear$^*$ & -1.9617 & 1.8890 & 0.0489 & -30.98 & 0.0684\\
\bottomrule
\end{tabular}}
\end{table*}

The adopted F150W relation is
\begin{equation}
 \Mbar_{150,0}=-3.208\pm0.024_{\rm stat}\pm0.051_{\rm common}
 \quad{\rm mag},
 \label{eq:f150}
\end{equation}
with intrinsic scatter 0.081 mag. The statistical error comes from
2000 galaxy-level bootstrap zero points. The alternative linear slope
0.2315 has an approximate $1\sigma$-equivalent bootstrap interval
$[-0.2977,0.7217]$, including zero. It worsens AICc by 3.14 and the LOO
RMS from 0.0925 to 0.0981 mag. Colour is therefore not required for the
combined F150W measurements in this sample.

For F090W, the adopted calibration is
\begin{equation}
 \begin{split}
 \Mbar_{090,0}={}&(-1.9617\pm0.0157_{\rm boot})\\
 &+(1.8890\pm0.4424_{\rm boot})(C_0-0.5703),
 \end{split}
 \label{eq:f090}
\end{equation}
with intrinsic scatter 0.0489 mag. Its linear colour term improves AICc
by 9.15 and reduces the LOO RMS by 0.0308 mag.
A logarithmic relation changes AICc by only 0.33 and LOO RMS by
0.0011 mag relative to the linear model, insufficient reason to adopt
a less direct population law. F150W nonlinear experiments likewise
do not improve both model-selection criteria materially.

\begin{figure*}[t]
\centering
\paperfigure{materials/pasa/short/calibration.pdf}{\textwidth}
\caption{Joint calibration comparison. The solid black line is the
adopted constant in F150W (left) and linear colour relation in F090W
(right); the dashed line is the alternative primary model.
Each grey strip shows only the adopted $\pm\sigma_{\rm int}$, not
a confidence interval on the fitted mean or a full distance error.
Red squares denote Fornax, blue circles the broad Virgo region, and
grey triangles other environments; labels are NGC numbers. All 14
galaxies enter each fit. Bars are individual $1\sigma$ colour and
absolute-SBF errors, excluding common zero points.}
\label{fig:calibrations}
\end{figure*}

The paired measurement variances in Figure~\ref{fig:measurement-variance}
show how the same image-error prescription contributes in each band.
They exclude the population and distance-calibration terms.

\begin{figure*}[t]
\centering
\paperfigure{materials/pasa/short/measurement_variance.pdf}{\textwidth}
\caption{Paired apparent-SBF measurement budgets. Each horizontal
segment is a contribution $\sigma_j^2$; the total includes foreground
extinction in addition to the internal terms of
Equation~\ref{eq:measurement-error}. Both panels use the same galaxy
order. These are not full distance uncertainties.}
\label{fig:measurement-variance}
\end{figure*}

\subsection{Comparison on the same galaxies}
\label{sec:same8}

For the eight objects common to this work and Paper III, the F150W
constant and linear models give LOO RMS 0.1014 and 0.1103 mag.
For Paper III F110W against $g-z$, a linear relation reduces RMS
from 0.1114 to 0.0465 mag. Changing the galaxy list therefore does
not explain the different colour response.

Figure~\ref{fig:same8} crosses three SBF bands with both available
colour indices. F150W has Pearson $r=0.28$ against F090W$-$F150W
and $r=0.21$ against $g-z$ ($p=0.503$ and 0.611).
F110W gives $r=0.82$ and 0.92 ($p=0.0127$ and 0.0010).
Its relation survives either colour choice, unlike the combined F150W
measurements. Passband is therefore an important difference, but
population response, distance prescription and reduction choices are
not isolated well enough to establish a unique physical explanation.

\begin{figure*}[t]
\centering
\paperfigure{materials/pasa/short/same8_colour.pdf}{\textwidth}
\caption{Same-eight-galaxy comparison across SBF bands (rows) and colour
indices (columns). Symbols identify environment as in
Figure~\ref{fig:calibrations}; all points carry NGC labels and individual
error bars. The ordinary least-squares lines are diagnostic, not the adopted
calibrations. Each single strip is $\pm s_{\rm res}$, with
$s_{\rm res}^2=\sum_i(y_i-\widehat y_i)^2/(8-2)$: observed residual
scatter including measurement noise, not fitted intrinsic scatter or
coefficient uncertainty.}
\label{fig:same8}
\end{figure*}

\subsection{Distance recovery and error budgets}

The adopted F150W and F090W relations give LOO RMS 0.0925 and 0.0684 mag
and biases $-0.0010$ and $-0.0039$ mag, respectively
(Figure~\ref{fig:recovery}).
Median internal predictive errors are 0.0865 and 0.0587 mag; including
the common TRGB scale gives 0.0984 and 0.0752 mag.
Validation-pull RMS values are 1.08 and 1.12. All F150W targets and
13 of 14 F090W targets lie within $2\sigma$; NGC 4649 is the
$2.49\sigma$ F090W exception and is retained.

\begin{figure*}[t]
\centering
\paperfigure{materials/pasa/short/distance_recovery.pdf}{\textwidth}
\caption{LOO distance recovery for the adopted F150W constant (left)
and F090W linear calibration (right). Each target is omitted from its
calibration; dashed lines mark equality. Horizontal bars are individual
TRGB errors and vertical bars are internal SBF predictive errors.
The common absolute scale is excluded. Environment symbols and NGC
labels follow Figure~\ref{fig:calibrations}. This is internal
cross-validation, not an independent external test.}
\label{fig:recovery}
\end{figure*}

Figures~\ref{fig:measurement-variance} and \ref{fig:distance-variance}
separate the error in measuring a fluctuation from the error in using
it as a distance indicator. Stacked segments add \emph{variances}, in
mag$^2$, not standard deviations. Small image errors do not imply
equally small distance errors: fitted population scatter and finite
calibration dominate much of the predictive budget.



\begin{figure*}[t]
\centering
\paperfigure{materials/pasa/short/distance_variance.pdf}{\textwidth}
\caption{Paired full LOO distance-error budgets. Segments show the
individual measurement, colour/reddening propagation, intrinsic scatter,
finite calibration and shared TRGB-scale variances. Colour contributes
only for F090W. The common scale is displayed inside each total but
is correlated between galaxies, not an independent per-object error.}
\label{fig:distance-variance}
\end{figure*}

\subsection{Sensitivity and remaining measurement limitations}
\label{sec:limitations}

Paper IV thick-slice and Paper-III-style cluster anchors give F150W
constant-model LOO RMS of approximately 0.093 mag, close to the primary
result; cluster tests separate depth from shared mean-distance errors.
Fixed quality cuts do not justify removing objects or restore a useful
F150W colour law. The Virgo--Fornax step remains diagnostic, not a detection.

Radial behaviour is not identical. F150W inner and outer annuli have
slopes 1.01 and 0.11 and intrinsic scatters 0.050 and 0.090 mag.
F090W slopes are 2.62 and 1.18. Both rings are retained to match the
published angular geometry, not because they are independent or equally
stable. A null combined F150W colour result is not a claim of no local
population gradient.

Internal PSF-size tests give maximum 129--257-pixel shifts of 0.0167 mag
in F150W and 0.0128 mag in F090W; an empirical isolated-star validation
is still missing. The F150W $0$--$2P_r$ test changes weighted magnitudes
by at most 0.0053 mag, and the special 1.35-mag globular-cluster widths
for NGC 1399 and NGC 4472 change them by at most 0.00047 mag.
This bounds the tested correction, not catalogue incompleteness itself.

Relative to the former raw-residual winsorisation, the adopted
normalised-space F150W branch shifts the combined magnitude by a median
$-0.02954$ mag and up to $0.08798$ mag in absolute value. The order of
operations is therefore a material methodological choice, not a negligible
correction.

The adopted noise model is imperfect. Median spectral reduced
$\chi^2$ values are 19.94 (F150W) and 4.48 (F090W), even after conjugate
mode counting is corrected. All adopted F150W ring fits have $P_1<0$;
the intercept approximates a correlated-noise continuum locally and is
not negative physical noise power. An empirical $N(k)$ extension
is nearly degenerate with $E(k)$ (correlations 0.982--0.997) and
does not provide a stable alternative. Archived synthetic FFT and
winsorisation checks constrain normalisation, not full reduction
systematics; they predate the bin-error correction and do not validate
its uncertainty coverage. Detailed diagnostics are in the Technical notes.

\FloatBarrier
\twocolumn[{
\begin{minipage}{\textwidth}
\section{Adopted distances}
\label{sec:distances}

Table~\ref{tab:distances} is the final distance comparison. Every error
is a full absolute $1\sigma$ uncertainty in Mpc. TRGB IV and both JWST
columns include the common 0.047-mag TRGB scale. The Paper III F110W
reconstruction includes measurement and colour errors, slope uncertainty,
0.060-mag population scatter and its 0.063-mag common scale; its
0.016-mag transfer term is already included in that scale.
The separate Jensen et al. (2015) F160W reconstruction includes
measurement, colour and coefficient errors, 0.114-mag scatter and the
0.10-mag common Cepheid zero point \citep{Jensen2015,Jensen2021,Jensen2025}.
Common terms must be cancelled or represented through covariance when
comparing columns; the printed errors are appropriate to quoting each
absolute distance in isolation.

\par\medskip
\begin{minipage}{\textwidth}
\tableshowtrue\centering
\captionof{table}{Final distances in Mpc with full absolute $1\sigma$
errors. TRGB IV: individual thin-slice anchors; Jensen III: HST F110W;
Jensen 2015: HST F160W. This work: colour-independent F150W and
colour-dependent F090W LOO distances. Ellipses indicate absent comparison
measurements. Shared scale terms remain correlated.}
\label{tab:distances}
\papertable{\textwidth}{\input{../../runs/F150W/analysis/tables/go3055_final_distance_comparison.tex}}
\end{minipage}
\par\medskip
\end{minipage}
}]

\section{Discussion}

F090W provides the tighter internal distance calibration on these
14 galaxies. F150W remains competitive without colour, while a common
population law for both filters is not supported. The comparison with
the same eight Paper III galaxies shows that the difference is not
simply a change in sample membership or choice of colour index.
However, the significant inner-annulus F150W trend and the limited
colour range prevent a universal assertion that F150W is colour independent.

Under the full-error convention, median uncertainties are 0.70 Mpc
for the eight Paper III F110W overlaps, 0.75 Mpc for F150W and
0.57 Mpc for F090W. These medians describe different overlap sets
and distances; they are not a controlled demonstration of universal
instrumental superiority. The five F160W comparisons give an
SBF-modulus offset $-0.119\pm0.133$ mag using joint covariance,
which does not establish a significant absolute-scale mismatch.
The Technical notes give the conditional covariance assumptions.

LOO tests common reductions and a shared anchor scale, so it cannot
measure external accuracy or uncover an offset common to every
calibrator. Residual structure, correlated noise, empirical PSF
normalisation and formal colour errors remain important limitations.
No internal-dust or object-specific SBF $K$-correction has been applied.
All targets have $z<0.007$, but an ordinary integrated-light
$K$ estimate is not a bound on the luminosity-moment-weighted SBF
correction \citep{Cervino2008}. No empirical internal-dust upper limit
is claimed. Cosmological flux dimming is already represented through
luminosity distance; an extra Tolman term is not added \citep{Hogg1999}.
An independent sample and more complete noise/population treatment
are needed before claiming external precision at the internal LOO level.

\section{Conclusions}

The same image-measurement procedure yields two useful but distinct
TRGB-anchored calibrations. All 14 targets are retained.
A constant $\Mbar_{150,0}=-3.208$ mag gives F150W LOO RMS
0.0925 mag; colour does not improve that prediction.
F090W instead favours a linear colour slope $1.889\pm0.442$ and
reaches 0.0684-mag LOO RMS. The median full absolute errors are
0.0984 and 0.0752 mag, respectively. These are internal results with
explicit common-scale and unresolved measurement systematics,
not an independent validation of the absolute distance scale.

\section{Data and code availability}

The calibrated JWST observations are public in MAST under GO-3055,
\href{https://doi.org/10.17909/z9ch-wk24}{doi:10.17909/z9ch-wk24}.
Derived tables, image products, analysis code and the accompanying
Technical notes are retained in the version-controlled project.
The implementation uses NumPy, SciPy, Astropy, Photutils,
Matplotlib and STPSF \citep{Harris2020,Virtanen2020,Astropy2022,
Hunter2007,Perrin2014}; exact processing versions and input-header
contexts are recorded in the Technical notes. The 28 input mosaics
have mixed archive pipeline/CRDS contexts rather than a uniform local
reprocessing. A permanent public code/data release has not yet been
assigned; this statement does not claim a local repository is public.

% Before submission: confirm authors/contact, funding/competing interests,
% acknowledgements, public archive/DOI, and the required AI-use disclosure.
% No scientific or administrative declarations have been invented here.
\FloatBarrier

\begin{thebibliography}{99}
\bibitem[Astropy Collaboration et al.(2022)]{Astropy2022}
Astropy Collaboration, Price-Whelan, A. M., Lim, P. L., et al. 2022,
ApJ, 935, 167, doi:10.3847/1538-4357/ac7c74.

\bibitem[Cervi\~no et al.(2008)]{Cervino2008}
Cervi\~no, M., Luridiana, V., \& Jamet, L. 2008,
\emph{On surface brightness fluctuations: probabilistic and statistical
bases. I. Stellar population and theoretical SBF},
arXiv:0809.4491, doi:10.1051/0004-6361:20077515.

\bibitem[Chazov et al.(2026)]{Chazov2026}
Chazov, M. I., Makarov, D. I., Tully, R. B., et al. 2026,
\emph{The TRGB--SBF Project. IV. A Color Calibration of the TRGB in the
JWST F090W+F150W Filters}, arXiv:2603.11160v1.

\bibitem[Harris et al.(2020)]{Harris2020}
Harris, C. R., Millman, K. J., van der Walt, S. J., et al. 2020,
Nature, 585, 357, doi:10.1038/s41586-020-2649-2.

\bibitem[Hogg(1999)]{Hogg1999}
Hogg, D. W. 1999, \emph{Distance measures in cosmology},
arXiv:astro-ph/9905116.

\bibitem[Hunter(2007)]{Hunter2007}
Hunter, J. D. 2007, Computing in Science \& Engineering, 9, 90,
doi:10.1109/MCSE.2007.55.

\bibitem[Jensen et al.(2003)]{Jensen2003}
Jensen, J. B., Tonry, J. L., Barris, B. J., et al. 2003,
ApJ, 583, 712, doi:10.1086/345430.

\bibitem[Jensen et al.(2015)]{Jensen2015}
Jensen, J. B., Blakeslee, J. P., Ma, C.-P., et al. 2015,
ApJ, 808, 91, arXiv:1505.00400.

\bibitem[Jensen et al.(2021)]{Jensen2021}
Jensen, J. B., Blakeslee, J. P., Ma, C.-P., et al. 2021,
ApJS, 255, 21, doi:10.3847/1538-4365/ac01e7.

\bibitem[Jensen et al.(2025)]{Jensen2025}
Jensen, J. B., Blakeslee, J. P., Cantiello, M., et al. 2025,
\emph{The TRGB--SBF Project. III. Refining the HST Surface Brightness
Fluctuation Distance Scale Calibration with JWST},
ApJ, 987, 87, doi:10.3847/1538-4357/addfd6.

\bibitem[Perrin et al.(2014)]{Perrin2014}
Perrin, M. D., Sivaramakrishnan, A., Lajoie, C.-P., et al. 2014,
Proc. SPIE, 9143, 91433X, doi:10.1117/12.2056689.

\bibitem[Reid et al.(2019)]{Reid2019}
Reid, M. J., Pesce, D. W., \& Riess, A. G. 2019,
ApJL, 886, L27, doi:10.3847/2041-8213/ab552d.

\bibitem[Tonry \& Schneider(1988)]{TonrySchneider1988}
Tonry, J. L., \& Schneider, D. P. 1988, AJ, 96, 807.

\bibitem[Virtanen et al.(2020)]{Virtanen2020}
Virtanen, P., Gommers, R., Oliphant, T. E., et al. 2020,
Nature Methods, 17, 261, doi:10.1038/s41592-019-0686-2.

\end{thebibliography}
\label{LastArticlePage}
\end{document}
